% Transcription par lecture visuelle directe, sans OCR.
% Statut : transcrit ; relecture indépendante à effectuer.
% Chapitre complet : p. 145 (fin)–166 (début).
% Source : PDF original, vue 072, côté droit, page manuscrite 145, après le chapitre quatrième.
\chapter{Première suite de la préfecture de Monsieur Chaulet}
\begin{center}
\textit{Révision du règlement. Travaux des sections, \&…}
\end{center}

Pendant que le cercle des relations de la Congrégation s'aggrandissait ainsi et réclamait de son gouvernement une attention continuelle pour les former et les entretenir, d'autres
% Source : PDF original, vue 073, côté gauche, page manuscrite 146.
devoirs exigeaient aussi tous ses soins. Il fallait que le préfet veillât sans cesse sur la Congrégation, sur l'observation de son règlement et en même temps fît prendre de l'activité à ses bonnes œuvres.

\section*{Révision du règlement}
% Mention marginale de la page manuscrite 146.
Parmi ces sollicitudes diverses, le préfet dut surtout songer à faire exécuter l'article du règlement qui fixait à dix-huit mois l'époque à laquelle ce règlement devait être révisé ; le moment en était arrivé ; il nomma donc une commission pour préparer ce travail ; les délibérations de cette commission se prolongèrent pendant plusieurs mois. Enfin elle soumit le règlement modifié à l'approbation de la Congrégation réunie en assemblée générale ; celle-ci y donna sa sanction après une légère délibération. On remarquait dans le nouveau règlement plutôt quelques additions importantes que des changements notables. La commission avait cherché à y introduire toutes les mesures que la Congrégation avait prises par arrêté depuis sa fondation ; ainsi une troisième section ajoutée au titre du gouvernement de la Congrégation (intitulé dans le nouveau règlement : du gouvernement de la Congrégation et de l'emploi de ses membres) établissait la division de la Congrégation en quatre sections ; quelques articles furent ajoutés au titre des réceptions pour indiquer la part que devait y prendre la section du zèle ; ici le règlement subit une modification assez importante ; jusqu'alors il avait suffi pour être admis de réunir en sa faveur la majorité des voix ; ce mode de recrutement fut trouvé beaucoup trop large ; c'était laisser une porte ouverte par laquelle des bataillons entiers auraient pu passer sans jamais rencontrer la moindre opposition.
% Source : PDF original, vue 073, côté droit, page manuscrite 147.
Ce mode avait pu être avantageux pour hâter la formation de la Congrégation dans un moment d'ailleurs où il offrait peu d'inconvénient, car tous ceux qui devaient alors faire naturellement partie de la Congrégation naissante étaient si bien connus des premiers membres, que leurs chefs pouvaient être considérés comme indiqués d'avance ; mais une fois la Congrégation constituée, il devenait urgent de changer cet état de choses ; en soumettant donc les sujets proposés à réunir à l'avenir la presque unanimité des suffrages, on établit la même proportion suivant le nombre des votants, qui subsiste encore.

On fit entrer dans le nombre des fêtes que nous devions célébrer en commun et auxquelles les réceptions devaient avoir lieu, la fête d'actions de grâces, déjà fixée au 4\up{e} dimanche après Pâques par un arrêté, et on en sortit la fête de la Pentecôte par rapport aux inconvénients que nous avons déjà remarqués.

Le dernier article détermina le mode qu'on devrait employer à l'avenir dans le cas qu'une révision ou des modifications à ce règlement fussent jugées nécessaires.

\section*{Sections}
% Mention marginale de la page manuscrite 147.
Les légères modifications qu'avaient subies le règlement n'en apportaient aucunes dans l'esprit qui dirigeait la Congrégation ni dans la marche qu'elle suivait ; cette nouvelle charte ne donna donc lieu à aucune secousse ; les sections commencèrent à prendre de l'activité et à mettre de la régularité dans leurs travaux ; celle du zèle donna la première l'exemple de se créer un règlement et de se choisir un patron ; elle se mit sous la protection de St François de Sales ; jusques là ses assemblées n'étaient pas tenues régulièrement ; elle les fixa à tous les huit jours ; elle introduisit l'usage d'une courte méditation
% Source : PDF original, vue 074, côté gauche, page manuscrite 148.
à chaque séance ; tous ces exemples furent bientôt suivis par les autres sections qui jusques alors ne s'étaient également assemblées que sur la convocation de leurs chefs ; la section de l'instruction prit pour patron St Augustin ; celle des aumônes St Jean de Dieu ; celle des prisons et hôpitaux St Vincent de Paul ; pour conserver en tout le principe de l'unité le préfet exigea que tous les règlements particuliers des sections fussent soumis à son approbation.

\section*{Section du zèle mise hors ligne}
% Mention marginale de la page manuscrite 148.
Il est dans la destinée des conceptions humaines de laisser toujours quelque chose à désirer ; il faut que partout l'imperfection de notre nature se laisse apercevoir ; la Congrégation n'est pas exempte de cette loi commune à laquelle l'humanité tout entière est assujettie ; à peine le nouveau règlement était-il en vigueur qu'on s'aperçut qu'en laissant la section du zèle former comme cela avait eu lieu jusques là, une fraction de la Congrégation, on avait fait une véritable faute ; la nature de son travail ne pouvait pas fournir un aliment suffisant à la bonne volonté de ses membres et cela d'autant moins que l'importance de ses fonctions exigeait qu'on la composât des meilleurs membres de la Congrégation ; c'était donc autant de bons ouvriers dont le zèle était paralysé et dont cependant le concours eût été bien utile dans les autres sections. Les termes du règlement se trouvèrent heureusement assez vagues pour qu'ils pussent laisser un vaste champ à l'interprétation ; le préfet n'hésita pas à en user ; il annonça qu'à l'avenir les membres qu'il désignerait pour faire partie de la section du zèle seraient en même temps attachés à une des trois autres sections. Seulement pour alléger un peu le surcroît
% Source : PDF original, vue 074, côté droit, page manuscrite 149.
de travail qu'on leur imposait, on régla que la section du zèle ne s'assemblerait plus que tous les quinze jours, ce qui pour elle était bien suffisant. Ce petit coup d'état eut un plein succès ; il rendit à la Congrégation toutes ses forces disponibles, et elle en profita pour redoubler de zèle dans la pratique des bonnes œuvres.

\section*{Section de l'instruction}
% Mention marginale de la page manuscrite 149.
Jetons maintenant un coup d'œil sur les progrès qu'ont pu faire les autres sections dont les débuts que nous avons vus au commencement de cette histoire, vous ont peut-être fait sourire de pitié. Dès l'origine, la section de l'instruction avait commencé l'exercice de ses fonctions en préparant quelques jeunes gens à la première communion ; bientôt on lui confia l'instruction de quelques personnes d'un âge mûr. Plusieurs de MM. les curés de la ville imaginèrent d'employer la bonne volonté de ses membres pour établir dans leurs paroisses des cours d'instruction, sous le nom de catéchismes de persévérance, en faveur des jeunes gens qui venaient de faire leur première communion. C'était un moyen bien avantageux pour fortifier ces jeunes gens dans la connaissance des vérités de la religion et les maintenir dans leurs bonnes dispositions. La section accepta très volontiers l'exercice d'une œuvre qui rentrait si directement dans ses attributions. Plusieurs de ces catéchismes eurent beaucoup de succès ; on peut signaler en particulier ceux de St Georges et de St Bonaventure. Le dernier surtout prit en peu de temps un accroissement très remarquable ; on y comptait plus de cinquante jeunes gens, parmi lesquels plusieurs âgés de 18 à 21 ans. Cette section joignait encore à ces catéchismes
% Source : PDF original, vue 075, côté gauche, page manuscrite 150.
qui étaient ses œuvres habituelles les œuvres de conversion qui pouvaient convenablement être tentées par elles ; les œuvres de ce genre ne pouvaient être opérées que difficilement par des jeunes gens qui jusqu'ici avaient plus de zèle que d'expérience. Elles ont donc été en petit nombre ; néanmoins la section réussit, dès ces premiers temps, dans la conversion de quelques protestants ; elle aurait voulu essayer aussi celle d'un Juif, mais sur les observations de M. Cholleton, vicaire général (c'était l'oncle du père actuel de la Congrégation), cette œuvre fut abandonnée. Dans le cours de mon long ministère, nous a dit ce vénérable ecclésiastique, j'ai vu quelquefois des protestants se convertir sincèrement mais pour des Juifs, jamais. La Congrégation n'a eu que trop, plus tard, par sa propre expérience, la confirmation de cette triste sentence. La section de l'instruction se chargea aussi quelquefois dans ces premiers temps des instructions à faire dans les prisons, mais le préfet trouva que cela établissait un peu de confusion entre nos sections et décida que la section qui était chargée de tout le service des prisons, devait également y faire les instructions.

\section*{Section des aumônes}
% Mention marginale de la page manuscrite 150.
L'œuvre de la section des aumônes ne présentait pas de grandes difficultés, aussi cette section ne fut pas bien longtemps avant d'avoir passablement grossi la liste de ses pauvres. Tout comme à présent, elle les visitait, leur portait secours et les encourageait à la pratique de la vertu ; tout comme à présent elle cherchait à glisser l'aumône spirituelle sous l'aumône temporelle afin de la faire parvenir plus sûrement à son adresse. Elle considérait toutefois la distribution des secours
% Source : PDF original, vue 075, côté droit, page manuscrite 151.
temporels comme sa fonction principale. M. Ducoin fut appelé à présider cette section en 1804. Le bien qu'il y a fait est incalculable ; l'excellent effet que sa présence y a produit s'y est fait longtemps sentir. Il introduisit dans cette section des habitudes d'ordre et de régularité qui facilitèrent beaucoup ses opérations. Ce fut lui qui commença à y faire tenir un registre où se trouvaient inscrits les noms de tous les pauvres adoptés par la section ; à côté de chaque nom se trouvaient aussi les motifs de leur admission et tous les renseignements qu'on pouvait se procurer sur leurs besoins, leurs ressources et leur moralité. Une fois par an à l'entrée et à la sortie de l'hiver, une longue séance extraordinaire (on y consacrait ordinairement toutes les heures d'un jour de dimanche que les offices laissaient libres) était employée à faire de ce registre un examen sévère, et par suite à débattre contradictoirement et à nouveau, comme si tout était en question pour la première fois, les secours qu'on accorderait à l'avenir à chaque pauvre. Les membres de la section n'étaient chargés personnellement d'aucun pauvre, mais ils étaient indistinctement et successivement visités par eux tous. Ces visites se faisaient toujours à deux ; elles avaient lieu pour chaque pauvre tous les 15 jours. La liste générale des pauvres était fractionnée en deux divisions, sous le nom de liste n\up{o} 1 et liste n\up{o} 2. À chaque assemblée de la section, on partageait tour à tour l'une de ces deux listes, en autant de petites listes particulières que le nombre des pauvres d'une part et celui des membres de la section de l'autre pouvaient l'exiger, afin que tout le monde fût employé et que personne ne fût surchargé. Ces listes étaient distribuées aux membres accolés de deux en deux ;
% Source : PDF original, vue 076, côté gauche, page manuscrite 152.
dans la semaine ils visitaient les pauvres dont ils avaient été chargés, et rendaient compte de leurs visites à la séance suivante. Les secours de la section des aumônes se distribuaient alors principalement dans le quartier de St Georges. Nos jeunes gens s'y étaient fait connaître par les bienfaits dont ils étaient porteurs, aussi ils y étaient en vénération ; on ne savait pas leurs noms, on mettait un grand soin à les laisser ignorer, mais on les honorait ouvertement du beau nom des messieurs de la miséricorde.

\section*{La Martin}
% Mention marginale de la page manuscrite 152.
Parmi les pauvres que visitait alors cette section d'aumônes qui dans les premiers temps de la Congrégation a brillé d'un si vif éclat, il y avait un personnage qui pour notre société est tout à fait historique ; son souvenir doit donc trouver ici sa place ; je suis convaincu qu'un grand nombre d'entre vous en a au moins entendu parler ; quant aux congréganistes des temps anciens, en lisant ce qui précède, ils ont déjà prononcé le nom révéré de la femme Martin.

La section des aumônes était presque à son début ; une pauvre femme lui est recommandée ; deux de ses membres sont délégués pour savoir ce dont il s'agit. Ils trouvent dans une petite chambre de la rue Tramassac, une pauvre malheureuse déjà avancée en âge, qui depuis 18 ans n'avait pas quitté, ni le jour ni la nuit, le fauteuil délabré sur lequel ses infirmités l'avaient reléguée. D'horribles ulcères couvraient ses jambes devenues aussi grosses que son corps ; ses plaies étaient affreuses à voir, elles étaient dans un état de suppuration continuelle, exhalaient une odeur de putridité, et lui faisaient éprouver des douleurs intolérables ; cette pauvre affligée était la femme Martin. Elle
% Source : PDF original, vue 076, côté droit, page manuscrite 153.
supportait sa fâcheuse position, sa misère qui en était l'inévitable suite, et les souffrances qu'elle endurait, avec une patience admirable. La section fut si édifiée des rares vertus qu'elle avait rencontrées dans ce triste réduit, qu'elle crut devoir déroger en faveur de la pauvre infirme, à la règle sagement établie de ne jamais visiter les maisons où il n'y avait point d'homme ; mais encore cette œuvre fut considérée en quelque sorte comme devant être l'œuvre fondamentale de la section. On peut dire que la section des aumônes devint dès lors, disons mieux, que la Congrégation tout entière est devenue comme la famille de la femme Martin. Le bruit de ses vertus se répandit parmi nous ; chaque congréganiste ambitionna le bonheur d'en être témoin, un grand nombre fut lui rendre visite. À la vue de tant de patience et de tant de résignation, on était si émerveillé, qu'on éprouvait aussitôt un désir bien vif de jouir de nouveau de cet édifiant spectacle. Ces visites étaient pour la vénérable infirme, un sujet de grandes consolations. Cette pauvre chambre était souvent un point de rendez-vous pour les congréganistes, il n'était pas rare de voir un cercle nombreux autour de ce fauteuil de douleurs. Lorsque MM. de Montmorency et de Noailles se trouvaient passer ou séjourner à Lyon, ils y paraissaient toujours, et souvent plusieurs fois. C'était ordinairement presque le premier mot que nous adressait en arrivant M. de Montmorency ; en réglant avec nous l'emploi de son temps pendant son séjour, il ne manquait guère d'ajouter : mais nous irons bien voir la Martin. Rien n'était plus touchant que cette réunion pieuse dont une pauvre infirme était le centre ; la conversation de ceux qui se rencontraient dans ce lieu si méprisé du monde et en même temps l'objet des complaisances du Seigneur, était dans le ciel ; on s'entretenait
% Source : PDF original, vue 077, côté gauche, page manuscrite 154.
des miséricordes de notre bon maître, des douceurs qu'on goûte à son service, et en portant les yeux sur ce modèle de patience dont on n'eût pas soupçonné le triste état en lui voyant le sourire sur les lèvres, on se sentait pénétré d'admiration et de reconnaissance envers le Dieu de bonté qui sait ainsi adoucir par sa grâce ce qui paraît si pénible et si rebutant à notre pauvre nature. La femme Martin est enfin allée recevoir la récompense de sa résignation à la s\up{te} volonté de Dieu. Comme elle en avait été entourée pendant sa vie, elle a encore vu les congréganistes autour d'elle à son heure dernière ; ils ont récité auprès d'elle les prières des agonisants et ont reçu son dernier soupir ; si, comme tout peut nous le faire présumer, elle a quitté cette terre pour aller chercher ce bonheur qui a été promis à ceux qui sont pauvres et à ceux qui pleurent, elle peut certainement être considérée comme une des protectrices de la Congrégation.

\section*{Section des prisons et hôpitaux}
% Mention marginale de la page manuscrite 154 ; placée au début du passage.
Ainsi que nous l'avons déjà dit, la section des prisons et hôpitaux avait de suite commencé à visiter quelques malades à l'Hôtel-Dieu. Tout le monde étant libre de se présenter pour cet objet à l'hôpital, aucune difficulté ne lui avait empêché d'accomplir cette partie de sa mission, et quoique cette œuvre n'eût pas encore obtenu l'extension qu'elle a acquise depuis, néanmoins elle se développait graduellement et commençait à prendre de la consistance. Ce n'était encore là qu'une branche des travaux dont devait s'occuper cette section ; il fallait aussi qu'elle parvînt à se faire ouvrir les prisons. Toute la première année s'écoula sans qu'elle pût trouver moyen d'y réussir ; mais avec la grâce de
% Source : PDF original, vue 077, côté droit, page manuscrite 155.
Dieu et une volonté bien prononcée, on finit presque toujours par triompher des obstacles. La prison qui paraissait réclamer davantage l'exercice de son zèle était celle de St Joseph ; c'était la prison des condamnés ; les détenus y étaient en grand nombre. Il vint heureusement à la connaissance des membres de cette section qu'un vénérable ecclésiastique, nommé M. Martinet, s'était entièrement dévoué au soulagement des malheureux que renfermait cette maison ; il avait obtenu la permission d'y établir une chapelle, et y disait la messe tous les dimanches. Cette découverte fut un trait de lumière pour la section. Elle lui traça la marche qu'elle devait suivre ; elle députa donc deux de ses membres auprès de M. Martinet ; ils devaient exposer le désir de plusieurs jeunes gens de venir au secours des prisonniers, se mettre à ses ordres, et le supplier de leur faciliter l'entrée de la prison de St Joseph.

Non seulement M. Martinet les accueillit avec bonté, mais avant de les introduire dans la prison, il les mit parfaitement au courant de la situation de la maison de St Joseph. Il n'était pas seul à s'occuper du soulagement des prisonniers, depuis plusieurs années de pieuses filles s'étaient mises en possession de cette bonne œuvre.

Pour votre édification et aussi à cause des rapports qu'elle a eus avec la section des prisons, je dois vous raconter l'origine de la société que formèrent entre elles ces saintes filles.

La fondatrice était une pauvre domestique dont le souvenir ne s'est conservé que sous le nom qu'elle portait en qualité de chrétienne, elle s'appelait Pierrette. Pierrette donc animée de l'esprit de foi, et sans autres ressources que les légères aumônes qu'elle pouvait recueillir entreprit pendant la Révolution
% Source : PDF original, vue 078, côté gauche, page manuscrite 156.
de soulager les prêtres et les émigrés qui, victimes de la tyrannie du gouvernement de l'époque, languissaient dans les prisons. La sœur Charlotte vint la première s'associer à ses bons œuvres ; quelques bonnes filles s'unirent ensuite à elles ; dès lors elles commencèrent à se considérer comme formant une société et d'un commun accord elles reconnurent Pierrette leur fondatrice comme leur supérieure. Lorsque après la Révolution qui plaça Bonaparte à la tête des affaires publiques, les prisonniers politiques manquèrent à leur zèle, ces pieuses filles ne voulurent pas pour cela discontinuer la bonne œuvre ; elles étendirent leurs soins sur tous les prisonniers ; elles se chargèrent de leur apporter des aliments, ce qu'elles firent aussi souvent que purent le permettre les aumônes qu'elles avaient recueillies. Au moment de la distribution, Pierrette faisait former un cercle par tous les prisonniers, se plaçait au centre, se mettait à genoux, leur intimait l'ordre d'en faire autant, faisait un grand signe de croix et prononçait une prière que tous répétaient avec elle. Elle leur apprenait ainsi à remercier celui qui donne la nourriture aux enfants des hommes tout comme aux petits des oiseaux. La mort seule mit fin aux bonnes œuvres de Pierrette. La sœur Charlotte lui succéda en qualité de supérieure ; bientôt l'œuvre s'aggrandit considérablement, on prit insensiblement l'habitude de l'appeler, œuvre des Charlottes, du nom de la seconde supérieure ; le nom lui en est resté. Des dames charitables établirent plus tard une société pour protéger cette œuvre qui subsiste toujours, elle continue à porter la soupe dans les prisons, surtout à la prison militaire dont elle est exclusivement chargée ;
% Source : PDF original, vue 078, côté droit, page manuscrite 157.
et quoiqu'elle soit aujourd'hui réunie aux sœurs de St Joseph, elle ne conserve pas moins son ancien nom d'œuvre des Charlottes, qui rappelle sa modeste et sainte origine.

Je vais même joindre ici le récit d'une petite circonstance de la vie de Pierrette, qui montrerait au besoin combien elle était vénérée par les malheureux qu'elle assistait. Un jour je ne sais par quelle affaire elle allait modestement à pied de Vienne en Dauphiné. C'était dans les dernières années de la Révolution. Dans ces temps de désordre, la route n'était rien moins que sûre ; Pierrette venait de traverser la petite ville de St Symphorien d'Ozon quand tout à coup elle se voit entourée par une bande de voleurs. Sans autre préambule, on lui demande sans façon la bourse ou la vie. Que peut-elle riposter à ce fâcheux compliment ? Sa bourse ? Elle n'en a pas guère, mais sa vie !… Elle en a encore besoin pour la consacrer à l'exercice de la charité. Elle se mettait donc en devoir de la racheter avec sa pauvre petite bourse, lorsque le chef de la bande s'approche, la reconnaît et s'écrie : misérables ! que faites-vous là ! C'est notre mère. À l'instant elle reçoit de tous les voleurs les témoignages du respect le plus profond ; ils s'empressent de se procurer un âne (St Symphorien, comme vous ne l'ignorez pas, fut jadis célèbre par sa poste aux ânes). Ils font monter Pierrette sur le pacifique animal, forment son escorte, et l'accompagnent ainsi en triomphe jusqu'à la proximité de Vienne ; là elle n'avait plus à craindre de fâcheuses rencontres.

Revenons à notre section des prisons et hôpitaux. Son concours venait fort à propos pour le soulagement des prisonniers, car ils ne recevaient guère d'autres secours que les aumônes que M. Martinet pouvait recueillir en leur faveur, et les distributions de soupe que
% Source : PDF original, vue 079, côté gauche, page manuscrite 158.
les faibles ressources des Charlottes ne leur permettaient pas de donner plus de deux fois par semaine.

Qu'il fut affreux le spectacle que présenta alors à nos congréganistes l'intérieur de la prison ! On en jugerait bien mal si on voulait prendre pour point de comparaison l'état où se trouvent ces maisons de détention, aujourd'hui que depuis bien des années les soins paternels d'une administration bienfaisante y ont établi des travaux et ont amélioré si notablement le sort des malheureux qu'elles renferment. Dès que l'énorme clef du geôlier eut fait, tour à tour, crier deux serrures, et que la porte de fer eut tourné péniblement sur ses gonds, pour leur ouvrir l'entrée de cet horrible séjour, ils furent presque suffoqués par l'odeur fétide qui s'exhalait de toutes parts. Ils aperçoivent des êtres, à figure blafardes et décharnées, ressemblant plutôt à des cadavres qu'à des créatures vivantes. Ces misérables, dévorés par la vermine, sont à demi couverts par de vieux pantalons délabrés et quelques lambeaux de chemises ; ils sont couchés sur une paille humide et à moitié pourrie ; nourris d'un pain noir que le gouvernement ne leur accorde qu'avec parcimonie ; et pour tout adoucissement à une affreuse position, les aumônes sollicitées par un St prêtre en leur faveur, et les foibles suppléments de nourriture que peuvent leur apporter quelques pauvres filles. Voilà jusqu'à présent, tout ce que la charité a pu leur offrir.

Aussi les cœurs des membres de la section furent-ils vivement émus. Ils se hâtèrent de se mettre en mesure pour concourir de tout leur pouvoir au soulagement de ces infortunés ; ils commencèrent par leur faire quelques distributions de vêtements ; peu à peu ils parvinrent à vêtir tous les prisonniers ; bientôt ils purent
% Source : PDF original, vue 079, côté droit, page manuscrite 159.
leur procurer la jouissance d'avoir enfin du linge propre et de pouvoir ainsi échapper à l'ordure et à la malpropreté, dans laquelle ils croupissaient depuis si longtemps. Tout en soignant leurs corps, ils cherchaient aussi à seconder le zèle de leur vertueux aumônier pour le bien de leurs âmes. Tous les dimanches deux membres de la section assistaient à la messe et veillaient à l'ordre. Le soir ils revenaient chanter vêpres avec les prisonniers et faisaient suivre cet exercice d'une instruction familière sur les principaux devoirs du chrétien. En un mot l'œuvre des prisons qu'on avait eu tant de peine à établir, prit de suite les plus grands développements. Cette section devint sans conteste à cette époque la portion la plus active de la Congrégation.

Dès que la porte des prisons eut été ouverte au zèle des congréganistes, par notre installation à St Joseph, la Congrégation des dames se hâta d'y envoyer aussi sa section des prisons ; mais cette section ne se borna pas seulement à remplir auprès des prisonnières, la même mission dont la nôtre était chargée à l'égard des prisonniers ; elle vint encore à son secours dans ce qui était de sa compétence. Ces dames voulurent bien se charger de confectionner les chemises que nous destinions aux prisonniers, et de faire également les autres travaux d'aiguille que pouvaient exiger nos distributions de hardes. Elle rend donc ainsi un grand service en venant au secours de la caisse de notre section, car dans l'état où se trouvait la prison, où tout était alors à créer, il eût fallu pour subvenir à tout ce qui eût pu paraître indispensable, d'autres ressources que celles dont quelques jeunes gens réunis pouvaient disposer. On fit toutefois des efforts inouïs, on dépensa dans cette seule année plus de 1\,500 f. pour le
% Source : PDF original, vue 080, côté gauche, page manuscrite 160.
service des prisons, et on parvint à couvrir cette énorme dépense ; mais il fallut pour cela employer tous les moyens possibles pour se procurer de l'argent. Les imaginations des membres de la section ne restaient pas en repos ; ils trouvèrent le secret de frapper à propos le rocher pour en faire sortir une source d'eau vive, et d'employer l'argent de l'iniquité à se faire des amis qui pussent les recevoir un jour dans les tabernacles éternels. Depuis longtemps des masses considérables de fausse monnoie s'étaient accumulées dans toutes les caisses de la ville ; personne n'avait songé à en tirer parti. Celui qui a toujours conduit avec sa main paternelle notre Congrégation, inspira à un membre de la section des prisons cette bonne pensée ; on fit la quête des écus faux. Les négociants qu'il embarrassait ne demandèrent pas mieux que de le voir disparaître ; on le porta à l'affinage, deux fontes eurent lieu successivement ; la première a rendu environ 400 f. et la seconde 240. Les membres de la section inventèrent encore un autre moyen de se faire de l'argent. L'or n'était point alors comme aujourd'hui considéré comme marchandise ; si parfois le besoin de s'en procurer faisait accorder une légère prime sur les pièces d'or, les chefs de commerce qui n'étaient pas en usage d'en tirer bénéfice, en abandonnaient le produit à leurs caissiers. Or, il y avait dans la section plusieurs caissiers ; ils s'accordèrent à faire tourner ce petit profit à l'avantage des prisonniers. À ces moyens extraordinaires que le génie inventif de la charité avait su découvrir, on joignit ceux qu'on pouvait exercer prudemment dans le cœur de
% Source : PDF original, vue 080, côté droit, page manuscrite 161.
personnes charitables, en faveur des prisonniers, sans toutefois compromettre le secret de la Congrégation. Enfin on ajouta encore à ces ressources diverses, le produit d'une quête qui fut faite publiquement un jour de dimanche, dans l'église de la Charité, par les dames et par nous. M. Martinet servait de manteau pour voiler les deux sections qui partageaient le gâteau à l'ombre de ses ailes.

Ce bon M. Martinet éprouvait une grande joie en voyant tout le bien qui commençait à s'opérer au milieu de ces prisonniers dont il se considérait comme le père ; il chercha les moyens de le rendre durable. M. Fay de Sattonai était depuis peu maire de Lyon, il administrait notre ville avec une grande sagesse ; M. Martinet lui suggéra la pensée de former dans l'intérêt des prisonniers, une administration gratuite qui gouvernerait paternellement les prisons, de la même manière que l'administration des hospices gouverne les maisons dont le soin lui est confié. Il suffisait de montrer à M. de Sattonai quelque moyen nouveau d'opérer le bien pour qu'il s'empressât de le saisir. Cette administration fut donc formée et a servi de modèle aux commissions des prisons qui plus tard ont été établies dans toutes les prisons de France. M. Martinet avait pris tellement en amitié les membres de la section des prisons, qu'il aurait voulu en composer la nouvelle administration ; mais d'une part, nos congréganistes étaient pour lors beaucoup trop jeunes ; de l'autre le caractère administratif devait imprimer un respect qui pouvait diminuer la confiance, ce qui aurait nui très certainement à l'œuvre que la section était appelée à remplir ; aussi on laissa chercher ailleurs les membres du conseil des prisons, et la section resta dans son incognito. Seulement elle fit bientôt entrer la prison de Roanne dans ses attributions et ce
% Source : PDF original, vue 081, côté gauche, page manuscrite 162.
ne fut que plus tard qu'elle s'introduisit aux recluses, où est la prison militaire.

\section*{Commission centrale des Congrégations}
% Mention marginale de la page manuscrite 162.
Ce n'était point seulement aux prisons que la Congrégation des dames et celle des jeunes gens cherchaient à s'entraider ainsi que vous venez de le voir. Ces deux Congrégations étaient toujours disposées à se soutenir et à se prêter main forte, toutes les fois que cela pouvait être utile aux bonnes œuvres. Il en était de même de la Congrégation des hommes, et même de celle des demoiselles, quoique la réserve qui fait son essence ne lui permît pas de donner à son concours une action aussi directe. Toutes ces Congrégations avaient le P. Roger pour père ; c'étaient donc comme autant de branches de la même famille.

Pour mettre de l'ensemble dans leurs bonnes œuvres on jugea convenable de former entre les Congrégations, une espèce de commission centrale. À cause de l'esprit de réserve dont nous avons déjà parlé, celle des demoiselles ne crut pas devoir y prendre part. Les trois autres envoyèrent chacune deux ou trois députés à ce comité central qui se réunissait tous les mois. On s'assemblait chez une des dames de la Congrégation. Là on arrêtait ensemble toutes les mesures qui pouvaient être utiles aux œuvres mixtes ; on compulsait les listes des pauvres, et s'il se trouvait quelque double emploi, on prenait les moyens d'y remédier.

La Congrégation paraissait se consolider tous les jours davantage, sa marche acquérait de la fermeté, ses œuvres prenaient de l'extension, tous ses membres étaient remplis de zèle, et cependant au milieu de cette prospérité réelle, un très petit incident donna lieu à une manifestation, véritable symptôme d'un mal qui pouvait y produire une révolution funeste. Qui croirait, Messieurs, qu'en
% Source : PDF original, vue 081, côté droit, page manuscrite 163.
parcourant les fastes de la Congrégation, on dût y rencontrer cet esprit d'insubordination et de révolte, qui se trouve à toutes les pages de toutes les histoires des peuples de tous les temps et de tous les lieux ? Cependant rien n'est plus certain, tant il est vrai que notre pauvre nature humaine se montre toujours partout avec ses misères. Toutefois rassurez-vous, nos émeutes n'ont point été sanglantes ; nos insurgés n'ont pas élevé de barricades ; ils n'ont point renversé de dynastie ; il a suffi d'un peu de fermeté de la part du préfet pour tout faire rentrer dans l'ordre.

\section*{Assistants isolés des sections}
% Mention marginale des pages manuscrites 163–164.
Voici le fait : la section des aumônes, établie en pleine possession de tous les avantages que lui avait procuré l'excellente administration de M. Ducoin, ses membres habitués à jouir aussi en commun du bonheur de marcher ensemble dans la carrière des bonnes œuvres, s'étaient fortement attachés les uns aux autres. Le jour de la fête patronale de la section arriva ; sa célébration fut accompagnée d'une petite agape ; là les cœurs se dilatèrent et s'épanchèrent à l'envi ; dans l'ivresse de la joie, dont le principe sans doute était bien pur, mais dont les effets furent bien désordonnés, ils se promirent réciproquement de ne jamais se séparer, lors même que le préfet viendrait à les désigner pour faire partie d'une autre section ; dans ce cas tout en se rendant à leur nouvelle section pour obéir, ils n'en continueraient pas moins à assister aux séances de la section des aumônes et à prendre part à ses travaux. Dès que le préfet eut connaissance de ce petit serment du jeu de paume en diminutif, il entrevit tout le péril qui pouvait menacer la Congrégation par suite d'une détermination aussi inconsidérée, ce n'était rien moins qu'une petite république qui entreprenait de s'établir au milieu de la grande ; il comprit qu'il importait pour arrêter le mal, d'agir promptement et de déployer de la vigueur. Quatre
% Source : PDF original, vue 082, côté gauche, page manuscrite 164.
membres de la section quasi insurgée, furent à l'instant placés dans d'autres sections et remplacés par d'autres. Tout le monde obéit sans sourciller ; personne n'osa dire un seul mot, il ne fut plus question des engagements pris et tout continua à marcher comme auparavant. À cette mesure on en joignit une autre. Ce qui venait de se passer démontrait plus que jamais la nécessité d'avoir l'œil constamment ouvert sur les sections ; d'ailleurs l'administration du préfet était souvent contrariée par suite du placement des assistants dans les sections, car cet état de choses avait continué à subsister jusqu'à cette époque. Assez ordinairement les assistants étaient en même temps chefs de section ; il résultait de là qu'au lieu d'être, comme cela eût été à désirer, les conseillers du préfet, ils remplissaient auprès de lui les fonctions d'avocats défenseurs de leurs sections respectives ; de plus ne voyant que la section à laquelle ils appartenaient et n'ayant aucun rapport avec les autres, ils ne pouvaient connaître que très imparfaitement l'ensemble de la Congrégation. Toutes ces réflexions décidèrent le préfet à profiter de la circonstance pour opérer une réforme qui lui paraissait indispensable. Il déclara à la Congrégation qu'il entendait à l'avenir employer les assistants exclusivement à l'aider dans le gouvernement de la société ; qu'en conséquence ils n'appartiendraient plus à aucune section, mais qu'ils les visiteraient toutes concurremment avec lui. Ce notable changement compléta l'organisation de la Congrégation, telle qu'elle existe jusqu'à ce jour.

Vous, Messieurs, qui depuis si longtemps en appréciez maintenant tous les avantages, vous serez peut-être étonnés d'apprendre que cette
% Source : PDF original, vue 082, côté droit, page manuscrite 165.
mesure dont l'expérience a démontré l'incontestable utilité fut alors accueillie avec beaucoup de froideur et même avec une sorte de peine par presque tous les congréganistes, mais surtout par les membres de la section des aumônes. C'étaient eux qui avaient donné lieu par leur prétention inconséquente à s'isoler ; ils en furent bien punis, car ils se virent privés d'un président qui leur était cher. M. Ducoin cessa ses fonctions de chef de la section des aumônes, et demeura assistant. Dans les premiers temps qui suivirent ce petit coup d'état, on redoutait dans les sections d'y voir arriver les assistants, on les considérait comme les espions du gouvernement de la Congrégation. On s'observait en leur présence ; on cherchait même à leur dérober la connaissance des affaires de la section. Cependant peu à peu les esprits se sont calmés, les préventions se sont dissipées, et bientôt on les vit paraître, non seulement sans la moindre peine, mais encore tout comme aujourd'hui avec une grande satisfaction.

\section*{Section extraordinaire}
% Mention marginale de la page manuscrite 165.
Pour compléter ce qui a trait à l'administration intérieure de la Congrégation par M. Chaulet, je dois vous parler d'une mesure à laquelle la Congrégation se détermina alors et qui ne s'est pas renouvellée depuis. Une cessation d'ouvrage vint enlever presque entièrement à notre population ouvrière les ressources sur lesquelles reposent ses moyens d'existence. Toutes les sympathies des citoyens de Lyon furent fortement excitées par le détresse toujours croissante des ouvriers. La Congrégation mue par les sentiments de charité qui l'animent, ne voulut pas rester étrangère à ce mouvement. Elle constitua une section à laquelle on donna le nom de section extraordinaire. Cette section fut chargée d'exercer en faveur des pauvres ouvriers, les mêmes œuvres que la section des aumônes exerçait envers les autres pauvres. Pour composer cette section
% Source : PDF original, vue 083, côté gauche, page manuscrite 166, avant le chapitre sixième.
on choisit quelques membres de chacune des trois sections ordinaires auxquelles ils demeurèrent toujours attachés. La section extraordinaire n'avait été établie que pour exister temporairement. Aussitôt que la cause qui avait donné lieu à sa création eut cessé, on la supprima et en terminant ses fonctions, elle confia aux soins de la section d'aumônes le petit nombre de personnes secourues qu'on ne pouvait pas encore abandonner, et lui remit le petit excédant qui se trouva dans sa caisse.

Nous nous sommes occupés dans ce chapitre de nos détails d'intérieur, pour ainsi dire de nos affaires de ménage. Dans le suivant nous allons considérer la Congrégation sur un plus grand théâtre.
